% Generated draft: requires mathematical and source review.
\documentclass[11pt]{article}
\usepackage[margin=25mm]{geometry}
\usepackage{fontspec}
\setmainfont{Latin Modern Roman}
\usepackage{amsmath,amssymb,mathtools,unicode-math}
\setmathfont{Latin Modern Math}
\usepackage{enumitem,needspace,xurl,hyperref}
\setlength{\parindent}{0pt}
\setlength{\parskip}{5pt}
\setlength{\emergencystretch}{3em}
\newcommand{\N}{\mathbb{N}}
\newcommand{\Z}{\mathbb{Z}}
\newcommand{\Q}{\mathbb{Q}}
\newcommand{\R}{\mathbb{R}}
\newcommand{\C}{\mathbb{C}}
\newcommand{\F}{\mathbb{F}}
\newcommand{\A}{\mathbb{A}}
\newcommand{\PP}{\mathbb{P}}
\newcommand{\E}{\mathbb{E}}
\newcommand{\eps}{\varepsilon}
\newcommand{\dd}{\,\mathrm{d}}
\begin{document}
\section*{UnsolvedMath ID: 30002226}
\textbf{Source problem code:} catalogue entry
\subsection{Definitions and mathematical statement}
Let \(\mathcal Z\) denote the Jiang--Su algebra. A complex \(C^*\)-algebra \(A\) is \(\mathcal Z\)-stable, or absorbs \(\mathcal Z\), if
\[
A\cong A\otimes_{\min}\mathcal Z
\]
as \(C^*\)-algebras. Under the usual convention, an elementary \(C^*\)-algebra is one isomorphic to \(\mathcal K(H)\), the algebra of compact operators on a nonzero complex Hilbert space \(H\); nonelementary means not of this form. The supplied record does not explicitly specify this convention.

The phrase ``counterexample to the Elliott classification conjecture'' does not specify a unique class of algebras in the supplied record. Write \(\mathscr C\) provisionally for the intended class of counterexample algebras. With that unresolved parameter, the proposed assertion is
\[
\forall A\in\mathscr C,\qquad
\bigl(A\text{ is nonelementary}\bigr)
\Longrightarrow
A\not\cong A\otimes_{\min}\mathcal Z.
\]
For example, if counterexample means membership in a pair violating classification by a specified invariant \(\operatorname{Ell}\) on a specified class \(\mathscr D\), one would first need to define
\[
\mathscr C=
\left\{A\in\mathscr D:
\exists B\in\mathscr D,\quad
\operatorname{Ell}(A)\cong\operatorname{Ell}(B)
\ \text{and}\ A\not\cong B
\right\}.
\]
This is a possible specification of the missing parameter, not a choice supplied by the source. The record does not identify \(\mathscr D\), the precise invariant, the meaning of its isomorphisms, or whether this pair-based meaning of counterexample is intended.

To express the second formulation, fix a free ultrafilter \(\omega\) on \(\N\) and, for a unital \(C^*\)-algebra \(A\), use the norm ultrapower convention
\[
\ell^\infty(\N,A)=\{(a_n)_{n\in\N}:\sup_n\|a_n\|<\infty\},
\qquad
I_\omega=\{(a_n)\in\ell^\infty(\N,A):\lim_{n\to\omega}\|a_n\|=0\},
\]
\[
A_\omega=\ell^\infty(\N,A)/I_\omega,
\qquad
A_\omega\cap A'
=\{x\in A_\omega:xa=ax\text{ for every }a\in A\},
\]
where \(A\) is identified with its constant-sequence image. A unital embedding means an injective unital \(*\)-homomorphism. The source's second proposed assertion is
\[
\forall A\in\mathscr C,\qquad
\bigl(A\text{ is unital and nonelementary}\bigr)
\Longrightarrow
\nexists\text{ a unital embedding }
\mathcal Z\longrightarrow A_\omega\cap A'.
\]
The ultrapower convention and the choice and quantification of \(\omega\) are not specified in the record. The source calls this formulation equivalent to nonabsorption, but does not supply the hypotheses needed to justify that equivalence. In particular, separability is not stated, and the unital formulation does not itself address nonunital algebras. Both displayed assertions therefore remain provisional formulations of the supplied question, pending specification of the intended classification setting and applicable equivalence hypotheses. The supplied dated assessment describes the underlying report assertion as contextual information about known counterexamples and does not establish the universal claim.
\subsection{Short English statement}
The proposed question asks whether every nonelementary counterexample to the intended Elliott classification conjecture fails to absorb the Jiang--Su algebra \(\mathcal Z\). For unital algebras, the source also asks whether \(\mathcal Z\) fails to embed unitally into \(A_\omega\cap A'\). The class of counterexamples and the hypotheses supporting the claimed equivalence are unspecified. The supplied open label is dated curation, not an independently verified mathematical conclusion.
\subsection{Sources}
\begin{itemize}
\item Local registry: \texttt{data/unsolvedmath/problems.json}, with \texttt{catalog/problems.json} as the supplementary registry; supplied complete record with permanent UnsolvedMath ID \(30002226\), title ``Z-Stability of Elliott Classification Counterexamples,'' and exact source problem code \texttt{catalogue entry}. Collection: Oberwolfach Reports Open Problems. No individual proposer is supplied; the proposed year is recorded as 2012.
\item Supplied report citation: \emph{C*-Algebras, Dynamics, and Classification (2013)}, Oberwolfach Reports, pp.\ 3129--3209, DOI \texttt{10.4171/owr/2012/52}. \url{https://doi.org/10.4171/owr/2012/52}. The supplied citation labels the report 2013, while the DOI identifier and proposed year contain 2012; these details are retained without reconciliation.
\item Supplied literature source, labelled ``Oberwolfach report, original assertion'': \url{https://publications.mfo.de/bitstream/handle/mfo/3323/OWR_2012_52.pdf?isAllowed=y&sequence=1}.
\item Literature assessment dated exactly \texttt{2026-08-21}: the supplied status is ``open.'' The supplied assessment characterizes the record as a contextual assertion about known Elliott counterexamples and states that the universal non-\(\mathcal Z\)-stability claim is not established in this form. This is supplied dated curation, not independently verified literature assessment or source verification.
\item The registry labels the statement \texttt{corrected\_verified} and reports a manual semantic, typographic, and LaTeX review. Those supplied editorial labels do not resolve the missing mathematical hypotheses or constitute independent verification here.
\end{itemize}
\end{document}
